\begin{abstract}
Language-guided aerial navigation could help responders inspect damaged infrastructure, but disaster imagery creates a failure mode that conventional vision-and-language navigation does not directly address: the agent may reach a plausible region while remaining uncertain about whether the observed structure is actually damaged.
We introduce \ProjectName, a system and benchmark formulation for bird's-eye disaster vision-and-language navigation in georeferenced pre- and post-event imagery.
Its central mechanism is calibrated uncertainty-driven active reinspection.
A dual-temporal damage model produces calibrated class probabilities; when uncertainty is high and an action is expected to reduce it, the agent spends a limited budget to descend and recenter before committing to a damage judgment.
Building proposals come from an event-disjoint ResNet34 U-Net localizer adapted from the xView2 localization stage, with damage labels still produced by the dual-temporal classifier; hierarchical planning, a textual semantic map, and episodic memory support the loop.
None of these perception or planning modules is presented as a success-rate improvement.
Our current evidence is deliberately bounded.
Under a frozen event-disjoint split, temperature scaling reduces test ECE from 0.095 to 0.037 without changing accuracy, but increases holdout ECE from 0.046 to 0.084.
A difference-attention fusion matches concatenation accuracy on both partitions and does not recover the older holdout accuracy gain observed under a leaky split.
Replacing the weak YOLO proposer raises building localization quality (test Dice $0.750$, proposal recall@$0.5$ $0.524$) while leaving the scientific claims about calibration and reinspection unchanged.
On the same 40-task suite with three paired seeds, E1 (B0--B3) and E11
(six reinspection policies) both obtain zero strict success and produce no
damage-evidence observations or reinspection triggers.
Thus the localization gain does not transfer into end-to-end evidence
coverage in the evaluated stack.
We therefore treat calibrated active reinspection as an executable, auditable hypothesis rather than a demonstrated navigation or judgment improvement.
\end{abstract}
